\section{Lean formalization and verification scope}\label{sec:formalization}

The public Lean 4 development~\cite{repository} covers stages M1--M11,
including both three-row proofs and the principal extensions in this paper.
Its counts are defined from geometric tilings, rather than from the target
recurrences. A cell has zero-based coordinates in
\texttt{Fin m} $\times$ \texttt{Fin w}; a domino is an unordered pair of
adjacent cells, and a tiling is an exact cover. The definition
\texttt{Domino.a n} is the cardinality of the quotient of actual
$3\times2n$ tilings by the four coordinate symmetries.
The coefficient of $X^n$ in \texttt{orbitSeries} is $(a(n):\mathbb Q)$.

\subsection{Correspondence with the mathematical results}
Stages M1--M3 construct the board, coordinate action, terminal-hole
objects, tiling--path equivalences, and reflection and half-turn
completion maps with their inverse identities. Stages M4--M6 prove
the support and counting induction in Section~\ref{sec:half}, apply
Burnside to the actual quotient, and derive Theorem~\ref{thm:main}.
For $B=(1-X)(1-4X+X^2)(1-4X^2+X^4)$ and
$N=1-3X-4X^2+10X^3-4X^5+X^6$, they first prove $BA=N$.
Since $B(0)=1$, this gives $A=NB^{-1}$ in $\mathbb Q[[X]]$;
coefficient comparison gives the seventh-order recurrence and all seven
initial values. The main endpoint is
\path{Domino.three_row_main_theorem}.

The later stages extend this coverage as follows.
\begin{center}
\small
\begin{tabular}{cp{0.80\linewidth}}
\toprule
Stage & Correspondence with this paper\\
\midrule
M7 & Theorem~\ref{thm:stabilizers}: exact stabilizer classes of actual
tilings, restricted orbit quotients, orbit sizes, and divisibility.\\
M8 & Theorem~\ref{thm:odd}: general-height local coverings and paths,
the symmetry completion bijections, the odd-height horizontal-reflection
term, matrix generating functions, rationality, and an eventual recurrence.\\
M9 & Theorem~\ref{thm:five} and Appendix~\ref{app:certificate}: the five-row
reachable states, six-dimensional compression, terminal vectors,
polynomial certificate, rational functions, fourteenth-order recurrence,
and fourteen initial values.\\
M10 & Theorem~\ref{thm:asym}: the three-row closed form, the two uniform
sampling spaces, parity asymptotics, exponential error bounds, and
vanishing proportions of symmetric tilings and orbits.\\
M11 & Section~\ref{sec:blocks}: the full block classification, unique
decomposition at all empty cuts, geometric word equivalence, coordinate
symmetries, palindrome centers, and a second proof of the main series.\\
\bottomrule
\end{tabular}
\end{center}

Representative endpoints are listed below in stage order:
\begin{quote}\small
\path{Domino.three_row_stabilizer_counting}\\
\path{Domino.General.odd_height_main}\\
\path{Domino.FiveRows.five_row_main_theorem}\\
\path{Domino.Asymptotic.three_row_asymptotic_main}\\
\path{Domino.Blocks.three_row_block_main}.
\end{quote}
In the namespace \path{Domino.Asymptotic}, the parity limits are
also recorded as \path{probability_even_asymptotic} and
\path{probability_odd_asymptotic}.
Stage-specific declaration lists and statement probes in
\path{lean/docs/} and \path{lean/Domino/Audit/} expose the exact types
and their dependencies.

The block development proves
\path{Domino.Blocks.orbit_series_via_blocks} without importing the
half-rectangle counting modules. It shares the definitions, local
tiling--path correspondence, Burnside's lemma, and pure algebra with the
other route. This is independence of the global combinatorial arguments,
not independence of authorship or of all foundations.
The formal block parameter starts at zero: \texttt{arm false k} and
\texttt{arm true k} represent $A_{k+1}$ and $B_{k+1}$.
Empty boards, the one-row case, and the $k=1$ block boundary are covered.

\subsection{Pinned environment and audit}
We use the public source release \texttt{paper-en-v2.0}, at commit
\begin{center}
\small\texttt{a5c32a4b81cce150dbd0dd2398fdd9e90664a353}.
\end{center}
The release name refers to its archived manuscript; its Lean sources
already include M1--M11. The fixed environment is Lean 4.35.0 and
mathlib commit
\begin{center}
\small\texttt{56a005855ea17e0b0958f73ec801ae140b0a1679}.
\end{center}
The cumulative inventory has 83 mathematical modules and 1078 explicit
declarations, including 699 theorem declarations and auxiliary
definitions and instances. These are source-inventory counts, not 1078
separate results of the paper. The recorded cumulative acceptance
audits transitive axiom dependencies and permits only
\texttt{propext}, \texttt{Classical.choice}, and \texttt{Quot.sound}.
The audited sources must contain no placeholder proofs or extra
mathematical axioms; \texttt{native\_decide} is also excluded.

The acceptance also reruns the earlier stages and rebuilds project
sources in a separate directory without prior project build products,
sharing the checked, pinned third-party dependency cache.
It compares the resulting source hashes, theorem types, and cumulative
axiom output. This is a fresh build of the project, not a cache-free
rebuild of all of mathlib. The public records include successful local
acceptance and a successful GitHub Actions run for the pinned release.
Commands for reproducing the checks are in Appendix~\ref{app:repro}.

\subsection{Limits of the formal claim}
The formal action is given by the four explicit coordinate
transformations; a classification of all Euclidean isometries preserving
a rectangle is not formalized. At zero width the abstract group acts
trivially on the unique empty tiling. No claim is made about minimal
recurrence order, general-height asymptotics, even-height extensions,
or square-board symmetries beyond the statements proved here.
Kernel verification checks the formal statements and their proofs;
external mathematical and semantic review, including whether the formal
definitions capture every intended convention, remains outstanding.
It also does not establish bibliographic priority. The older
\path{review/m1-m6/} package retains its original, narrower review scope;
it is not an external review of M7--M11.
